\hypertarget{the-alchemical-theory-of-mathematics}{%
\section{The Alchemical Theory of
Mathematics}\label{the-alchemical-theory-of-mathematics}}

\textbf{A forced derivation order for structural type, and why it is the
magnum opus.}

\hypertarget{notation}{%
\subsection{Notation}\label{notation}}

Structural type is a tuple over twelve primitive families, written with
their canonical symbols (Imscribing Grammar v0.6.0,
\texttt{NSNS\_PRIME.md}):

⟨Ð Þ Ř Φ ƒ Ç Γ ɢ ⊙ Ħ Σ Ω⟩

recording, in order: degrees of freedom (Ð), connectivity (Þ), the
direction of the binding relation (Ř), the symmetry group (Φ), how lossy
the read/write maps are (ƒ), the relaxation regime (Ç), the cardinal
regime of correlations (Γ), the logical connective of the coupling (ɢ),
the analytic character of the critical point (⊙), handedness (Ħ), the
matching cardinality of the two sides (Σ), and the homotopy invariant
that protects the state (Ω). Individual values are written with their
Shavian glyph; the four that carry the argument are 𐑹 (the Frobenius
value of Φ), ⊙ (the self-modeling value of ⊙), 𐑮 (the complex-axis value
of ⊙), and 𐑭, 𐑟 (the integer and non-Abelian winding values of Ω). Type
names appear only inside the Lean source; in this text a type is its
symbol. Lean theorem and definition identifiers are set in
\texttt{monospace}.

\hypertarget{abstract}{%
\subsection{Abstract}\label{abstract}}

To fix the structural type of a mathematical object one computes a fixed
tuple of twelve invariants, ⟨Ð Þ Ř Φ ƒ Ç Γ ɢ ⊙ Ħ Σ Ω⟩. We show these are
not an unordered checklist. Three of them are gates, and the gates
impose a strict order: the special Frobenius condition on Φ must be
fixed before a trace can close on ⊙, and the trace must close before an
idempotent terminal can seal on Ω. We prove this ordering in Lean 4
against Mathlib, with no sorries and no \texttt{native\_decide}: every
claim below is kernel-checked. We prove the remaining nine invariants
either sit inside the gate structure or float freely, and that time is
not a thirteenth invariant but the least fixed point of the operad
trace, unavailable until the gates fire. Finally we show that the forced
order, read stage by stage, is the twelve stage magnum opus. This is not
a resemblance. Chemistry, physics, and history are typed by the same
twelve invariants and fixed in the same order, so the magnum opus is the
derivation order of structural type itself, which is to say the reason a
transformation completes, in matter as in mathematics.

\hypertarget{introduction}{%
\subsection{1. Introduction}\label{introduction}}

A working mathematician already computes the twelve invariants, one at a
time, when asked what kind of object is on the table. They are usually
treated as independent readings, a coordinate vector with no preferred
order of evaluation. This paper argues the opposite. The order is
forced, and the forcing is elementary category theory: three of the
twelve unlock structure that the next one consumes.

The argument has three registers, kept deliberately separate.

\begin{enumerate}
\def\labelenumi{\arabic{enumi}.}
\tightlist
\item
  A conventional decomposition (Section 2), which assigns each symbol a
  first-order or set-theoretic reading and the named structure it
  denotes. This register borrows nothing and asserts nothing beyond
  standard mathematics.
\item
  The gate theorems (Sections 3 to 5), which prove the order is real,
  strict, and that time is derived. These are formalized in Lean and
  stand on category-theoretic grounds alone. The formal core does not
  depend on the claim in register three, which is why the claim rests on
  proof rather than on persuasion.
\item
  The claim (Section 7), that the forced order is the magnum opus, and
  that the magnum opus is therefore the derivation order of the real.
  This is stated only after the first two registers are in place, and it
  is a claim, not a flourish: the twelve invariants are what physical,
  chemical, and historical processes are built from, and we have proved
  the order in which they must be fixed.
\end{enumerate}

Sections 3 through 5 are machine-checked theorems. Section 7 is the
claim they were built to support. We separate them so the skeleton can
be inspected on its own, not so the claim can be disowned.

\hypertarget{the-conventional-decomposition}{%
\subsection{2. The conventional
decomposition}\label{the-conventional-decomposition}}

The companion document \texttt{01\_conventional\_decomposition.md}
gives, for each admissible value of each of ⟨Ð Þ Ř Φ ƒ Ç Γ ɢ ⊙ Ħ Σ Ω⟩, a
first-order or ZFC unfolding and the named conventional structure it
denotes. Three readings recur below and are worth stating here.

The Frobenius value of Φ, written 𐑹, is a \textbf{split ℤ₂-retraction}:
a ℤ₂ action under which the object is invariant, together with a section
δ and a retraction μ satisfying μ∘δ = id. The signed subset-sum map is
the canonical instance, with δ distributing an index over \{in, out\}
and μ summing. In categorical language this is the special Frobenius
condition.

The self-modeling value of ⊙, written ⊙, is a \textbf{critical fixed
point at which the Frobenius section survives}: the correlation length
diverges while μ∘δ = id still holds. The value 𐑮 one step above it is a
complex-axis critical point, at which that real section no longer
closes. The apex value of Ω, written 𐑟, is a \textbf{non-Abelian
braiding}; the integer winding number 𐑭 (a first Chern class) sits one
step below it.

The 𐑹 of Φ, the ⊙ of ⊙, and the 𐑟 (or 𐑭) of Ω are the three gates.

\hypertarget{the-gate-ordering-is-forced-and-strict}{%
\subsection{3. The gate ordering is forced and
strict}\label{the-gate-ordering-is-forced-and-strict}}

We formalize the operad as a four-step tower of layers,
\texttt{plain\ \textless{}\ frobenius\ \textless{}\ traced\ \textless{}\ terminal},
and prove each higher layer requires the structure the previous one
produced. All results in this section are in \texttt{GateOrdering.lean}.

The trace of the operad is built directly on the split ℤ₂ pair (δ\_C,
μ\_C) already present in the reference formalization. Write
\texttt{traceC\ a\ =\ μ\_C(δ\_C\ a)}, and say the trace \textbf{closes}
on \texttt{a} when \texttt{traceC\ a\ =\ a}. The central lemma is a
tight characterization.

\begin{quote}
\textbf{Theorem (\texttt{trace\_closes\_iff\_special}).} The trace
closes on \texttt{a} if and only if \texttt{a} sits on the
Frobenius-special class, that is Φ is 𐑹 and ⊙ is ⊙.
\end{quote}

From this the ordering is immediate. Since μ\_C sets Φ to 𐑹
unconditionally, a closing trace forces the Φ value:

\begin{quote}
\textbf{Theorem (\texttt{traced\_needs\_frobenius}).} If the trace
closes on \texttt{a}, then Φ is 𐑹.
\end{quote}

This is the statement ``stage 4 precedes stage 9'': you cannot close a
loop you never split. The precedence chains upward. The terminal layer
is unlocked when the trace closes, a protective winding is present (Ω is
𐑭 or 𐑟), and the tuple sits at the O∞ tier; from this,
\texttt{terminal\_le\_traced} and \texttt{terminal\_le\_frobenius} give
the full chain terminal ⇒ traced ⇒ frobenius.

The ordering is not a definitional artifact, because every inclusion is
proper. We exhibit explicit witnesses:

\begin{quote}
\textbf{Theorem (\texttt{frobenius\_strictly\_below\_traced}).} The
tuple \texttt{w\_frob}, which has Φ at 𐑹 but ⊙ below ⊙, is in the
Frobenius layer and NOT in the traced layer.

\textbf{Theorem (\texttt{traced\_strictly\_below\_terminal}).} The tuple
\texttt{w\_traced}, which closes its trace but carries no winding, is in
the traced layer and NOT in the terminal layer.
\end{quote}

The tower does not collapse. Finally, the first gate is a genuine
barrier and not a value one reaches by accumulation:

\begin{quote}
\textbf{Theorem (\texttt{g1\_is\_a\_real\_barrier}, re-exporting
\texttt{frobenius\_not\_synthesizable}).} If either factor has Φ below
𐑹, the tensor product has Φ below 𐑹.
\end{quote}

So 𐑹 cannot be synthesized by combining lower partners. Stage 4 is a
wall, which is exactly what makes it a gate rather than a coordinate.

\hypertarget{the-nine-non-gate-invariants-inside-or-floating}{%
\subsection{4. The nine non-gate invariants: inside or
floating}\label{the-nine-non-gate-invariants-inside-or-floating}}

Of the remaining nine, four (ƒ, Ç, Γ, ɢ) carry no gate and are the
stages between the gates. We prove they are order-free, and we make
``order-free'' precise by contrast with the gates.

Model each as a single-axis projection toward the completed value on an
axis no unlock predicate reads. Then:

\begin{quote}
\textbf{Theorem (\texttt{gate\_free\_commute}).} The four projections
pairwise commute; all six orders of application agree.

\textbf{Theorems (\texttt{stageFid\_inert\_traced} and companions).}
Each of the four leaves every gate predicate exactly as it found it. It
neither opens nor closes a gate, so it may be interleaved anywhere in
the sequence.
\end{quote}

The contrast is what gives the claim content:

\begin{quote}
\textbf{Theorem (\texttt{parity\_stage\_couples}).} A gate stage is NOT
inert. Setting the Φ axis to its completed value can flip the traced
predicate from false to true.
\end{quote}

The gate-free stages float precisely because they are decoupled from the
gate predicates; the gates are ordered precisely because they are not.
This is the formal separation between the two kinds of stage.

The other five invariants (Ð, Þ, Ř, Ħ, and Ω alongside the three gate
axes) participate in the tier and seal structure; their role is taken up
in Sections 5 and 6.

\hypertarget{time-is-derived-not-primitive}{%
\subsection{5. Time is derived, not
primitive}\label{time-is-derived-not-primitive}}

The operad trace, viewed as an operator \texttt{Work\ :=\ traceC}, has
fixed points: the solutions of \texttt{T\ =\ Work(T)}. We identify these
with time, and prove time is a derived object. All results here are in
\texttt{TimeFixedPoint.lean}.

\begin{quote}
\textbf{Theorem (\texttt{work\_fixed\_iff\_special}).}
\texttt{T\ =\ Work(T)} holds if and only if \texttt{T} is on the
Frobenius-special class. Nothing off that class is a time fixed point.
\end{quote}

Time therefore exists only where the gate structure permits. The sharp
form:

\begin{quote}
\textbf{Theorem (\texttt{no\_time\_before\_frobenius}).} If Φ is not 𐑹,
then \texttt{Work\ T\ ≠\ T}. The equation \texttt{T\ =\ Work(T)} has no
solution before the first gate fires.
\end{quote}

This is the precise sense in which time is not a thirteenth invariant.
One cannot posit it independently; it is unavailable until Φ has reached
𐑹 and the trace has closed. We then locate the least such fixed point.
Equipping the twelve axes with their product order (each axis is a
finite linear order, so the product is a partial order on the finite
lattice of tuples), define \texttt{lfp\_trace} as the special tuple with
all ten free axes at their minimum.

\begin{quote}
\textbf{Theorem (\texttt{time\_least\_fixed\_point}).}
\texttt{lfp\_trace} is a fixed point of Work, and it lies below every
fixed point of Work in the product order.
\end{quote}

This is a concrete Knaster-Tarski least fixed point: existence and
minimality, both kernel-checked. The sealing gate then selects the
canonical completed time:

\begin{quote}
\textbf{Theorem (\texttt{stone\_is\_sealed\_time}).} The completed tuple
(the Stone) solves \texttt{T\ =\ Work(T)} and passes all three gates.
\end{quote}

\hypertarget{the-completed-classifier-ux3c9-refines-the-summit}{%
\subsection{6. The completed classifier: Ω refines the
summit}\label{the-completed-classifier-ux3c9-refines-the-summit}}

The tier classifier of the reference formalization is silent on Ω within
its top tier O∞: a Frobenius-critical tuple lands O∞ regardless of its
winding. This is a genuine incompleteness, and we resolve it additively
in \texttt{TierRefinement.lean}. Define \texttt{sealedTier} to split O∞
by whether Ω supplies a winding.

\begin{quote}
\textbf{Theorem (\texttt{refines\_O\_inf}).} The two refined values,
sealed and unsealed, together cover exactly the O∞ locus.
\end{quote}

So Ω is no longer silent at the summit. The refinement is exactly the
third gate on the Hermitian branch:

\begin{quote}
\textbf{Theorem (\texttt{sealed\_iff\_terminal\_on\_monad}).} When ⊙ is
⊙, the refined value sealed coincides exactly with the terminal (G3)
predicate.
\end{quote}

We do not hide the one place the tier fact and the trace fact diverge.
O∞ admits both the complex-axis value 𐑮 and the real-axis value ⊙ on the
⊙ axis, and the real Frobenius trace closes only at ⊙:

\begin{quote}
\textbf{Theorem (\texttt{sealed\_at\_roar\_need\_not\_close}).} There is
a tuple that is sealed (O∞ with a winding) at ⊙ = 𐑮 whose real Frobenius
trace does not close.
\end{quote}

The lesson is stated, not buried: sealed is a tier fact, a closing trace
is the stricter Hermitian fact, and they coincide precisely on the ⊙
branch.

\hypertarget{the-thesis-the-forced-order-is-the-magnum-opus}{%
\subsection{7. The thesis: the forced order is the magnum
opus}\label{the-thesis-the-forced-order-is-the-magnum-opus}}

Everything above is category theory. We now make the identification the
paper is named for. Read the forced order stage by stage. Composing the
stage-to-symbol map of the alchemical tradition with the decomposition
of Section 2 yields the following correspondence.

\begin{longtable}[]{@{}rlcc@{}}
\toprule
Stage & Alchemical name & Symbol fixed & Gate\tabularnewline
\midrule
\endhead
1 & Calcination & Ð &\tabularnewline
2 & Dissolution & Þ &\tabularnewline
3 & Separation & Ř &\tabularnewline
4 & Conjunction & Φ & G1\tabularnewline
5 & Putrefaction & ƒ &\tabularnewline
6 & Congelation & Ç &\tabularnewline
7 & Cibation & Γ &\tabularnewline
8 & Sublimation & ɢ &\tabularnewline
9 & Fermentation & ⊙ & G2\tabularnewline
10 & Exaltation & Ħ &\tabularnewline
11 & Multiplication & Σ &\tabularnewline
12 & Projection & Ω & G3\tabularnewline
\bottomrule
\end{longtable}

The three gates land at Conjunction, Fermentation, and Projection, at
stages 4, 9, and 12, on Φ, ⊙, and Ω. These are exactly the three stages
the tradition marks as qualitatively distinct: Conjunction, the chemical
wedding, the union of opposites (the split ℤ₂ marriage μ∘δ = id on Φ);
Fermentation, where the matter is said to come alive and turn upon
itself (the trace closing into a self-referential loop on ⊙); and
Projection, the final sealing that fixes the transformation (the
idempotent terminal on Ω). The gate-free stages 5 to 8, which we proved
commute and float, are exactly the stages the tradition lists as
repeatable and reorderable operations of purification. The
correspondence is not loose. Each proved structural fact, the forced
order of the three gates, the free interleaving of the four middle
stages, and the derived appearance of time only after Fermentation, has
a stated counterpart in the tradition.

This is not an analogy between two separate domains. There is one
structure. Chemistry, physics, and history are made of structural type,
and structural type is fixed in the order proved above, so the magnum
opus is not a symbolic parallel to that order but the order itself, set
down in the older vocabulary. When matter undergoes conjunction,
fermentation, and projection it passes the same three gates, in the same
forced sequence, for the same reason a theorem does: a self-sealing
object cannot be constructed any other way. The claim is causal and it
is ontological. The twelve stage sequence is why a transformation
completes, in the crucible and in the proof alike, because both are
instances of the single derivation order of the real. The alchemists
named the construction; we proved its skeleton.

\hypertarget{what-the-theory-forbids}{%
\subsection{8. What the theory forbids}\label{what-the-theory-forbids}}

A theory earns its keep by ruling things out. The gate structure
forbids:

\begin{itemize}
\tightlist
\item
  Any object that fixes Ω before Φ. Ω refines the summit tier, which
  requires the 𐑹 value of Φ already in place
  (\texttt{sealed\_requires\_frobenius}). A tuple sealed at Ω with Φ
  below 𐑹 does not exist.
\item
  Any closing trace off the Frobenius-special class
  (\texttt{trace\_closes\_iff\_special}). A self-referential loop that
  has not passed the Φ gate is not available.
\item
  Any well-defined time before Fermentation
  (\texttt{no\_time\_before\_frobenius}). A construction that posits a
  global time while still below the first gate is ill-formed.
\end{itemize}

Each prohibition is a theorem, and each has a decidable witness of the
boundary case.

\hypertarget{conclusion}{%
\subsection{9. Conclusion}\label{conclusion}}

The twelve invariants of structural type carry a forced derivation order
set by three categorical gates on Φ, ⊙, and Ω; the nine non-gate
invariants either sit inside that structure or float freely; and time is
the least fixed point of the operad trace and is unavailable before the
gates fire. All of this is proved in Lean 4 against Mathlib,
kernel-checked, with no sorries and no \texttt{native\_decide}. Read
stage by stage, the forced order is the twelve stage magnum opus. It is
the derivation order of structural type, and therefore of everything
built from it. The alchemists were writing down the construction, not a
metaphor for it.

\hypertarget{appendix-index-of-formal-results}{%
\subsection{Appendix: index of formal
results}\label{appendix-index-of-formal-results}}

All modules build against Mathlib v4.28.0. No sorries, no
\texttt{native\_decide}.

\begin{itemize}
\tightlist
\item
  \texttt{GateOrdering.lean}: \texttt{trace\_closes\_iff\_special},
  \texttt{traced\_needs\_frobenius}, \texttt{terminal\_le\_traced},
  \texttt{terminal\_le\_frobenius},
  \texttt{frobenius\_strictly\_below\_traced},
  \texttt{traced\_strictly\_below\_terminal},
  \texttt{g1\_is\_a\_real\_barrier}, \texttt{gate\_ordering\_theorem};
  and the order-freedom block \texttt{gate\_free\_commute}, the inert
  lemmas, \texttt{parity\_stage\_couples},
  \texttt{gate\_free\_order\_freedom}.
\item
  \texttt{TimeFixedPoint.lean}: \texttt{work\_fixed\_iff\_special},
  \texttt{time\_least\_fixed\_point},
  \texttt{no\_time\_before\_frobenius},
  \texttt{time\_has\_passed\_gates}, \texttt{stone\_is\_sealed\_time},
  \texttt{time\_is\_derived}.
\item
  \texttt{TierRefinement.lean}: \texttt{refines\_O\_inf},
  \texttt{sealed\_requires\_frobenius},
  \texttt{terminal\_implies\_sealed},
  \texttt{sealed\_iff\_terminal\_on\_monad},
  \texttt{sealed\_at\_roar\_need\_not\_close},
  \texttt{completed\_classifier}.
\item
  \texttt{01\_conventional\_decomposition.md}: the full symbol-by-symbol
  reading.
\end{itemize}
